This report summarises the current status of palimpsest and the results
of an initial benchmarking study, following the meeting of 19 June 2026.
It covers the two benchmarks agreed at that meeting, now executed on a
five-paper reference set; the expenditure to date; the design of the
knowledge graph and schema; and two interface components under active
development. The findings are presented as hypotheses to be validated at
larger scale, together with a costed proposal for that validation.
Unless noted otherwise, accuracy figures are computed against the
reference set as it stood when each run was executed: 41 measurements
for the parser-by-model grid (Section 3) and 47, the current reference
set, for the multi-pass study (Section 4); see Section 2.

\section{Status}\label{status}

The extraction and benchmarking infrastructure is complete, and the two
benchmarks proposed on 19 June have been executed on a five-paper
reference set: (i) a comparison of PDF parsers measured by downstream
extraction accuracy, with the parser as the independent variable, and
(ii) a study of language-model breadth spanning frontier, inexpensive,
and ensemble configurations. Both use a single deterministic scoring
rule and a common baseline design, so the two are directly comparable.
Two interface components have reached a minimum viable state: the agent
terminal interface and a browser-based verification viewer. The run-time
autonomy mechanisms discussed on 19 June have been implemented, namely a
schema-reading tool, a consistency gate linking skill, schema, and
ontology, run-time self-diagnosis of dropped extractions, and a domain
skill for the PEM water-electrolysis anode. The agent\textquotesingle s
bash escape hatch is now confined to the workspace by an
operating-system sandbox (macOS Seatbelt or Linux bwrap), which
strengthens the constrained-autonomy property that forms the core of the
thesis.

In summary, the infrastructure and the benchmark are complete on five
papers; scale-up validation, further testing, and writing remain.
Expenditure to date is approximately \$35 (Section 5).

\section{Benchmark methodology}\label{benchmark-methodology}

\subsection{Corpus and ground truth}\label{corpus-and-ground-truth}

The corpus comprises five papers on the oxygen-evolution reaction and
PEM water electrolysis, drawn from the Nature family of journals (Table
1). The reference set was not produced by exhaustive manual labelling.
For each paper, several independent readings were aggregated, namely the
text extracted by the four parsers together with a direct language-model
reading of the PDF (the \texttt{*.pdfread.md} anchor), and these were
reconciled to a consensus and spot-checked against the raw parsed text,
with a human review of a sample. A value is admitted to the reference
set only if it appears in at least one parser\textquotesingle s text;
values that occur only within a figure are recorded as a per-parser
coverage ceiling rather than as model failures.

Each reference entry is a bare
\texttt{(measurement-type,\ numeric-value)} pair. The catalyst and the
experimental conditions are recorded in the source tables but are not
used by the matching procedure, a deliberate simplification whose cost
is discussed in Section 2.2. The reference set was curated iteratively,
and the audit trail is retained:

\begin{itemize}
\tightlist
\item
  From 41 to 40: a review (\texttt{experiments/results/gold\_audit.md},
  21 June 2026) removed one mislabelled entry,
  \texttt{Stability\ 2.5\ h}, which denotes a measurement window rather
  than an endurance result.
\item
  From 40 to 47: a later completeness review added seven in-scope
  measurements (six PEM full-cell voltages and one 400 h stability
  value) that the original list had omitted. Each candidate was admitted
  only after agreement among at least three models and subsequent
  verification against the parsed text; agreement served as a filter,
  not as a verdict.
\end{itemize}

Each published figure is therefore scored against the reference set as
it stood when that run was executed: the parser-by-model grid (Section
3, Table 2) against the pre-audit 41-measurement set, and the multi-pass
and ensemble study (Section 4) against the current 47-measurement set.
The two differ by the seven measurements added in the completeness
review. The frontier models in the grid have not yet been re-scored on
the 47-measurement set, so Section 4 compares the inexpensive ensemble
against the grid's frontier figure only approximately; that re-score
forms part of the scaled run (Section 8).

\textbf{Table 1.} The five-paper reference set (post-audit reference, 40
measurements).

\begin{longtable}[]{@{}L{0.24}L{0.16}L{0.23}L{0.06}L{0.13}@{}}
\toprule\noalign{}
Paper & Journal (year) & System & Pages & Reference measurements \\
\midrule\noalign{}
\endhead
\bottomrule\noalign{}
\endlastfoot
s41467-022-35426-8 & Nat. Commun. (2022) & Ir-Co₃O₄ single atoms & 12 &
19 \\
s41467-023-40912-8 & Nat. Commun. (2023) & IrO₂@TaB₂ & 12 & 7 \\
s41467-025-63541-9 & Nat. Commun. (2025) & Ir/TiOₓ@Ti & 14 & 4 \\
s41565-025-02030-y & Nat. Nanotechnol. (2025) & RuₓIrOₓ (figure-only
Tafel) & 12 & 8 \\
s41929-024-01168-7 & Nat. Catal. (2024) & amorphous IrOₓ vs rutile & 13
& 2 \\
\end{longtable}

The total is 40 after the audit. The grid in Section 3 was scored on the
pre-audit 41-measurement set (one additional stability value in the last
row, later removed); the verified extension adds three, one, and three
measurements to rows 2, 3, and 4 respectively, giving the current
47-measurement set (19, 10, 5, 11, 2).

The baseline design is shared by both studies. A floor configuration is
defined as the least-expensive model with a single reference parser and
no domain skill, concretely deepseek-v4-flash on MinerU, the locked
runtime default. Each study then varies one factor away from this floor:
the parser (the parser study), the language model (the model study), or
the skill.

\subsection{Scoring}\label{scoring}

Matching is performed by deterministic code, not by a language model and
not by regular expressions (\texttt{experiments/ab\_extract.py}). A
prediction matches a reference value when the measurement type is
identical and the value lies within the looser of a 1 \% relative
tolerance or an absolute floor (±0.5 for values of magnitude at least
one, ±1×10⁻⁴ below one). The reported metrics are precision (true
positives over emitted values), recall (true positives over reference
values), the F1 score, and the cost per correct extraction. Before
matching, values are normalised deterministically: canonical units are
checked (tolerating paper-faithful spellings, for example s⁻¹ and 1/s),
a generous per-slot magnitude ceiling catches unconverted prefixes, and
a narrow milli-prefix unit re-derivation is applied.

It is important to state what the scorer does and does not measure. It
measures numeric recall and precision on type and value only. It is
deliberately context-blind: a value is counted as correct even if it was
extracted for the wrong catalyst or condition, because the reference set
omits that context. Semantic correctness is therefore not scored at
present. It is, however, enforced when a value is written to the graph,
through ontology typing, closed-shape SHACL validation, and mandatory
provenance (Section 6). Condition-aware matching remains a known gap
rather than a feature. The language model is used only for extraction
(one call per page); scoring and validation are entirely deterministic.

On the question of DeepEval, raised on 19 June: DeepEval is designed for
language-model-as-judge scoring of free text, whereas the present task
is numeric extraction, which a deterministic scorer answers without
additional cost, without run-to-run variance, and without the
circularity of using a language model to grade language models during a
model benchmark. A deterministic scorer is therefore proposed as the
primary metric, with DeepEval reserved, if used at all, for an optional
fuzzy cross-check.

Four parsers enter the grid: docling, MinerU, dots.ocr, and PaddleOCR.
Chandra is excluded because it emits no bounding-box geometry and
requires about 17 minutes per paper (\$0.105 per paper); it is reported
as excluded on these grounds rather than as a zero
intersection-over-union score.

\section{Parser-conditional extraction
accuracy}\label{parser-conditional-extraction-accuracy}

Across the five-paper set, the choice of parser affects accuracy as
strongly as the choice of model. With the model and prompt held fixed, a
single model\textquotesingle s micro-averaged recall varies by more than
30 percentage points across parsers, and replacing MinerU with docling
raised raw micro-recall for eight of the nine models tested. docling
produces roughly an order of magnitude more, and finer-grained, text
spans than MinerU\textquotesingle s coarse blocks, so values reported in
figures and tables become citable; the figure-only Tafel slopes in the
RuₓIrOₓ paper rise from a ceiling of four of eight on the text-only
parsers to eight of eight on docling. This is the central empirical
claim of the thesis, demonstrated here at small scale.

\textbf{Table 2.} Best model per parser (raw extraction, micro-recall,
models with complete five-paper coverage; 41-measurement pre-audit
reference). The partial sonnet-on-PaddleOCR row is excluded as not
comparable.

\begin{longtable}[]{@{}L{0.12}L{0.23}L{0.10}L{0.45}@{}}
\toprule\noalign{}
Parser & Best model & Recall & Inexpensive alternative \\
\midrule\noalign{}
\endhead
\bottomrule\noalign{}
\endlastfoot
MinerU & deepseek-v4-pro & 88 \% & gemini-flash-lite 83 \% \\
docling & sonnet-4.6 & 95 \% & gemini-flash-lite 93 \% (about 12×
cheaper) \\
dots.ocr & gemini-3.5-flash (free) & 88 \% & (none; \$0) \\
PaddleOCR & deepseek-v4-pro & 93 \% & gemini-3.5-flash 92 \% (\$0) \\
\end{longtable}

For recovery ceilings the parsers rank docling ≳ dots.ocr ≈ PaddleOCR
\textgreater{} MinerU. Once MinerU is set aside, the better models are
sufficiently robust across parsers that the choice of model has the
larger effect on accuracy.

\section{Frontier, inexpensive, and ensemble
configurations}\label{frontier-inexpensive-and-ensemble-configurations}

The most accurate single model across the full parser grid is
gemini-3.1-flash-lite, with a mean µ-F1 of 0.80, recall of 81 to 93 \%
depending on the parser, and a cost of \$0.0011 per correct extraction
(the grid of Section 3, scored on the 41-measurement reference).
Robustness across parsers is the decisive property for an agent that may
invoke any parser.

The multi-pass and ensemble study was run on docling and PaddleOCR and is
scored on the current 47-measurement reference, so its figures are
reproducible directly from the stamped results table
(\texttt{experiments/results/llm\_matrix\_t74\_2026-06-21.csv}). The
single-run recall of an inexpensive model is consequently lower here than
the corresponding entry in Table 2: gemini-lite recovers the same 38
values in both cases, which is 93 \% of the pre-audit 41-measurement set
but 81 \% of the extended 47, the difference being the seven added PEM and
stability measurements that a single inexpensive run misses.

\begin{longtable}[]{@{}L{0.40}L{0.16}L{0.34}@{}}
\toprule\noalign{}
Configuration & Recall (47-measurement) & Effect \\
\midrule\noalign{}
\endhead
\bottomrule\noalign{}
\endlastfoot
gemini-lite, single run, docling & 81 \% (38/47) & baseline for the
strategies below \\
gemini-lite with re-query, docling & 87 \% (41/47) & a within-model
second pass recovers three further values \\
Union of deepseek-flash and gemini-lite (one run each), docling & 100 \%
(47/47) & the two models' errors are decorrelated; the union is complete \\
The same union, PaddleOCR & 91 \% (43/47) & near-complete; the residual is
not all recoverable from PaddleOCR's text \\
\end{longtable}

The most effective configuration is the union of two inexpensive models,
each run once. Because the two models omit different values, that is,
their errors are decorrelated, the union recovers all 47 reference
measurements on docling at negligible marginal cost, the individual runs
having already been performed; the seven PEM and stability measurements
that defeat a single inexpensive run are exactly those the union
recovers. A within-model re-query is more limited: it raises gemini-lite
from 81 to 87 \% on docling but does not reach completeness. The best
frontier single run on the earlier 41-measurement grid (sonnet-4.6 on
docling, 95 \%, Table 2) is of the same order as the union; a
like-for-like re-score of the frontier models on the 47-measurement
reference has not yet been performed and forms part of the scaled run
(Section 8). The proposition that an ensemble of inexpensive models
reaches frontier-level recall at a fraction of the cost is therefore
supported on docling, and is stated as hypothesis H2 for validation at
larger scale.

Multi-pass strategies are not uniformly beneficial.

\begin{itemize}
\tightlist
\item
  A reason-then-format strategy, in the manner of CRANE, fails on the
  weakest inexpensive model: deepseek-flash falls from 94 \% (its plain
  single-run recall) to 40 \% at an 8k output limit, because the
  reasoning field exhausts the token budget and truncates the items, and
  recovers only to 66 \% at a 16k limit, still below plain extraction. It
  does not transfer to an extraction model constrained by output length.
\item
  A targeted re-query is two-edged: it raises the stronger inexpensive
  model (gemini-lite, 81 to 87 \% on docling) but lowers the weaker one
  (deepseek-flash, 94 to 89 \%); the strategy must be matched to the
  model.
\item
  The locked default, deepseek-v4-flash, is the weakest model overall on
  the grid (mean µ-F1 0.47) and too sensitive to the parser to serve as a
  default; its low nominal cost is offset by its misses. We propose
  replacing it with deepseek-v4-pro (a steady 88 \% or so on every
  parser) or with gemini-flash-lite.
\item
  The most expensive model returned little: opus-4.8 cost \$0.0338 per
  correct extraction, about 30 times the least-expensive viable models,
  yet scored below deepseek-pro and sonnet on the grid, and was removed
  after the first stage.
\end{itemize}

The recommended production configuration is docling with
gemini-3.1-flash-lite, the best single parser-and-model pairing on the
grid (Section 3), supplemented where completeness matters by a second
inexpensive model whose output is unioned with the first; where a zero
marginal cost is decisive, the free gemini-3.5-flash is a strong baseline
on docling and PaddleOCR.

\section{Cost to date}\label{cost-to-date}

\textbf{Table 3.} Expenditure (cost ledger \texttt{palimpsest.db} and
actual RunPod usage, as of 24 June 2026).

\begin{longtable}[]{@{}L{0.34}L{0.08}L{0.48}@{}}
\toprule\noalign{}
Category & \$ & Notes \\
\midrule\noalign{}
\endhead
\bottomrule\noalign{}
\endlastfoot
LLM total & 25.98 & 901 metered calls across the model, parser, and
multi-pass grid \\
of which OpenAI (GPT-5.4, -mini) & 11.69 & the most expensive tier,
mid-ranked accuracy \\
of which Anthropic (sonnet, haiku, opus) & 8.41 & includes the pruned
opus-4.8 (\$1.73) \\
of which Google (gemini-lite, -flash) & 3.32 & best value, used most
(210 and 49 calls) \\
of which DeepSeek (pro, flash) & 2.09 & flash heavily tested (276
calls), proposed for retirement \\
of which other (qwen, minimax) & 0.47 & dropped models, data retained \\
GPU / RunPod (actual) & about 9.40 & five-parser cloud verification (see
note) \\
Total & about 35.4 & remaining headroom about \$14.6 \\
\end{longtable}

The in-process cost meter records only the GPU expenditure it drives
directly (\$1.80 across 12 sessions). Actual RunPod expenditure is
approximately \$9.40; the difference is un-metered pod wall-time
(earlier runs, image download, idle, and setup), which by design lies
outside the in-process meter and is the operator\textquotesingle s
responsibility. The honest total is therefore approximately \$35.4.

The full grid of four parsers and roughly nine models, together with the
multi-pass arms, cost approximately \$26 in language-model calls. The
inexpensive models dominate the accuracy-per-dollar frontier, with a
spread of about 30 times between opus and gemini-lite per correct
extraction; the most expensive models contributed little additional
accuracy.

\section{Knowledge graph, schema, and
ontology}\label{knowledge-graph-schema-and-ontology}

Each extracted measurement becomes a single RDF node, typed by its
ontology class and carrying a value (\texttt{xsd:float}), a unit label
(string), and an optional confidence. It is linked to two further nodes:
a condition node (current density, electrode potential versus the
reversible hydrogen electrode, temperature, scan rate, electrolyte
family and pH, cell type, iR correction, and so on) and an evidence node
typed \texttt{prov:Entity} that carries the provenance treated as
non-negotiable, namely the paper SHA-256, the page, four bounding-box
coordinates, the parser name, and the source text. Run metadata (run
identifier, parse run, and the extraction-model tag) is written to a
per-run named graph, which keeps the data graph closed under SHACL. A
measurement that lacks complete provenance is refused rather than
inserted. A single overpotential measurement emits about 22 quads (five
for the measurement, nine for the evidence, four for the paper, three
for the condition, and one in the run graph).

The schema is defined once, in a single LinkML file
(\texttt{schema/palimpsest.yaml}), which is the source of truth. From it
are generated the Pydantic models, the JSON schema, the JSON-LD context,
and the SHACL shapes, so that one contract simultaneously validates
extraction, shapes the language model\textquotesingle s tool signature,
and enforces conformance of the RDF. The schema defines a measurement
hierarchy of eleven types (overpotential, Tafel slope, exchange current
density, charge-transfer coefficient, mass activity, turnover frequency,
ECSA, specific activity, stability, PEM cell voltage, and degradation
rate), together with the structural classes Paper, Condition,
Electrolyte, and Evidence. New slots are proposed in
\texttt{schema/exploratory.yaml} and reviewed before they enter the
closed main schema; they are never added silently.

Ontology alignment operates at two levels. In the triples themselves,
each measurement node is typed by an EMMO ECHO electrochemistry class
IRI where one exists (for example, overpotential, exchange current
density, and ECSA), or by a palimpsest-local IRI where EMMO provides no
term (Tafel slope, mass activity, the PEM metrics, and the
oxygen-evolution reaction). Provenance uses PROV-O, and bibliographic
metadata uses schema.org. At the schema level, as annotations rather
than emitted triples, QUDT units and H2KG (version 1.0.0) equivalences
are attached through \texttt{skos:closeMatch}. H2KG already defines
three of the four metrics that were originally local, so aligning to it
also aligns the work to a recognised DECODE and Helmholtz standard.
Values are at present stored as a flat number with a unit string, rather
than in the relational form that H2KG uses (measurement, hasProperty,
quantity value); bridging to that form is future work.

The working RDF store currently holds 14,637 quads across 60
extraction-run named graphs. This reflects every model-and-parser
combination inserted during benchmarking, rather than the unique
measurements of five papers. A recently added, append-only corrections
path allows a human or the agent to supersede a value with an immutable
audit chain; this is a sanctioned second write path to the graph that
also serves as feedback to the extractor.

\section{Status of the findings}\label{status-of-the-findings}

All results above rest on five papers and 40 reference measurements, and
are stated as hypotheses for validation at larger scale:

\begin{itemize}
\tightlist
\item
  H1: a parser-robust inexpensive model (gemini-flash-lite) matches
  frontier single-run accuracy at roughly one tenth of the cost.
\item
  H2: the union of two inexpensive models with decorrelated errors
  closes the gap to frontier accuracy at negligible marginal cost.
\item
  H3: docling is the best general-purpose parser for downstream
  extraction, and the union of parsers is the reliable means of
  recovering figure-only values.
\end{itemize}

A risk to generalisation should be stated plainly. The extraction guards
(the magnitude ceilings and the milli-prefix unit re-derivation) were
tuned against these five papers and required three rounds of review
before they stopped corrupting their own test spans, which is a sign of
over-fitting. The 20-to-25-paper run will be treated as the test of
generalisation, and a guard that misfires will be relaxed rather than
further patched. Precision carries a caveat in both directions: a
completeness review of 105 apparent false-positive groups found that
about 32 (roughly 30 \%) are real values that the reference set omits,
so precision is understated, while about 52 (roughly 50 \%) are likely
hallucinations and represent the genuine loss. Precision is thus neither
as high nor as low as a raw count would imply, and it is the axis that
most needs the larger corpus.

\section{Proposed next phase and cost
projection}\label{proposed-next-phase-and-cost-projection}

The next phase comprises two activities of very different cost.

The first is an application, or demonstration, run, corresponding to the
use case discussed on 19 June: 20 to 25 papers on hydrogen technology
passed through the chosen pipeline (docling with gemini-flash-lite),
driven end-to-end by the agent. Parsing here uses docling only, so GPU
cost is small and extraction costs about \$0.15; the total is
approximately \$1 to \$2, with a small margin for any failures. This run
fits within the current budget.

The second is a full re-benchmark at the scale of 20 to 25 papers, which
would validate hypotheses H1 to H3 and re-score on the 47-measurement
reference. It re-runs the grid of parsers, models, and multi-pass arms
at five times the present corpus size:

\begin{longtable}[]{@{}L{0.70}L{0.20}@{}}
\toprule\noalign{}
Item & Estimate \\
\midrule\noalign{}
\endhead
\bottomrule\noalign{}
\endlastfoot
Parse 20 to 25 papers across four or five parsers (GPU, actual rates) &
about \$30 to \$40 \\
Model grid (four parsers, about eight models, multi-pass), about 5× the
five-paper spend & about \$120 to \$150 \\
Ground-truth construction (multi-reading, verification, human
spot-check) & principally an effort cost \\
Total compute & about \$150 to \$200 \\
\end{longtable}

The parse estimate uses actual RunPod rates, about five times the
in-process meter (Section 5); an earlier estimate of \$6 had
under-counted pod overhead.

The full re-benchmark is what would convert the present hypotheses into
defensible thesis claims, and it requires roughly \$150 to \$200, which
exceeds the budget available to the project. We would therefore be
grateful to know whether Forschungszentrum Jülich is able to support the
scaled validation, whether as API or GPU credits or as an increased
budget of the order of \$150 to \$200. The estimate is itemised above so
that the figure can be assessed independently, and the demonstration run
can proceed within the current budget in any case.

\section{Agent interface and verification
viewer}\label{agent-interface-and-verification-viewer}

Both components are functional and at a minimum viable state. Figures
are attached: a still of the terminal interface, and stills of the
viewer. A screencast of the agent terminal interface in use is available
at \url{https://youtu.be/bWA27LkR4hQ}.

The agent terminal interface (theme "Scriptorium") is a conversational
front end to the constrained-autonomy agent. Replies stream as they are
produced; each tool invocation appears as a collapsible trace with its
elapsed time, with the supervision-relevant tools (\texttt{bash} and
\texttt{extract\_paper}) shown in full and the others summarised. A
footer shows a budget gauge (a live expenditure bar that moves from
sepia through amber to red across graduated warning thresholds), an
issue badge, and the true context size, alongside clickable links and a
copy-reply action. Fifteen cost-free slash commands are available,
covering expenditure (\texttt{/cost}, \texttt{/budget}), model routing
(\texttt{/use\ \textless{}role\textgreater{}\ \textless{}provider\textgreater{}}),
the session (\texttt{/resume}, \texttt{/clear}, \texttt{/export},
\texttt{/undo}), and provenance and audit (\texttt{/view} for the
viewer, \texttt{/git} for the workspace checkpoints, \texttt{/issues},
and \texttt{/review} for an agent summary of the session), together with
switchable themes (\texttt{/theme}).

\begin{figure}[H]
\centering
\includegraphics[width=\linewidth]{tui-traces.png}
\caption{The agent terminal interface (theme ``Scriptorium''): a scripted extraction
showing the collapsible tool trace with its elapsed time and insert count, the streamed
reply, and the budget-gauge footer with the active model roster.}
\end{figure}

The verification viewer is a two-pane page (FastAPI with PDF.js): the
source PDF on the left, and on the right a parser-by-model matrix
selector together with three tabs that follow the extraction pipeline.
The first tab, Parser, shows the raw text spans the parser recovered;
the second, Schema, shows the typed measurement cards derived from them
(slot, value, unit, confidence badge, page, source text, and an edit
action); the third, Gold, shows that cell scored against the reference
set (true positives over total, recall, precision, and the matched or
missed status of each reference value). The Gold tab reuses the
benchmark\textquotesingle s own matching procedure, so the numbers a
reviewer sees are identical to the benchmark\textquotesingle s.
Selecting any card or span highlights its source region on the PDF; that
region is now narrowed to the exact value text by means of the PDF text
layer, giving a precise highlight rather than a whole paragraph. Parsers
that report pixel coordinates are scaled to the page, and parsers
without geometry decline to draw a box rather than draw a wrong one. The
viewer also supports append-only corrections (by a human or the agent),
CSV export, and graceful degradation. The tabbed viewer is complete;
narrowing the stored geometry across all parsers is in progress.

\begin{figure}[H]
\centering
\includegraphics[width=\linewidth]{viewer-funnel.png}
\caption{The verification viewer, Schema tab: the source PDF on the left, and on the
right the parser-by-model matrix selector with the typed measurement cards. Each card
links to its source region on the page.}
\end{figure}

\begin{figure}[H]
\centering
\includegraphics[width=\linewidth]{viewer-gold.png}
\caption{The verification viewer, Gold tab: the selected cell scored against the
reference set (matched values over total, recall, and precision) using the benchmark's
own matching procedure, so that the viewer's figures coincide with the benchmark's.}
\end{figure}

Together these components serve the Accessible dimension of the FAIR
principles, and the tabbed viewer mirrors the thesis itself: it follows
the same parser-to-schema-to-reference path that the benchmark measures,
with every value traceable, by a single click, to the region of the page
from which it was read.

\section{Open questions}\label{open-questions}

The following points would benefit from the supervisor\textquotesingle s
guidance.

\begin{itemize}
\tightlist
\item
  Scoring metric. We propose a deterministic scorer as the primary
  metric and suggest that DeepEval, if used at all, serve only as a
  secondary fuzzy cross-check. Would this be acceptable?
\item
  Parser count. We propose reporting four parsers in the grid and
  documenting Chandra separately as excluded (no bounding-box geometry;
  prohibitive run time). Is this appropriate?
\item
  Default model. We would like to replace the locked deepseek-v4-flash
  default with either deepseek-v4-pro or gemini-3.1-flash-lite, and
  would welcome your view.
\item
  Funding. Would Forschungszentrum Jülich be able to support the scaled
  validation run (Section 8, second item; approximately \$150 to \$200)?
\item
  Ground truth at scale. What reference protocol would you recommend for
  the 20 to 25 papers: the multi-reading-and-verification approach used
  here; how many papers should receive a human-checked reference; and
  would meaning-aware (condition-context) scoring be worth adding, given
  that it is not scored at present (Section 2.2)?
\item
  Any further comments. We would be glad of any further thoughts on the
  approach, whether on the benchmark design, the inexpensive-model and
  ensemble direction, the constrained-autonomy framing, or the H2KG
  alignment.
\end{itemize}

The figures and numbers in this report trace to
\texttt{experiments/results/FINDINGS.md}, the \texttt{*.meta.json}
snapshots, \texttt{experiments/corpus\_manifest.csv},
\texttt{experiments/results/gold\_audit.md}, the \texttt{palimpsest.db}
cost ledger together with actual RunPod usage, the live \texttt{store/},
and \texttt{schema/palimpsest.yaml} with
\texttt{src/palimpsest/store.py}. The reproducible benchmark reports are
\texttt{reports/t72\_report.html} and \texttt{reports/t74\_report.html}.
